\section{Introducción}

Entre 2008 y 2025, el mercado laboral colombiano atravesó transformaciones importantes. Este informe analiza cómo evolucionaron los ingresos laborales por hora a lo largo de casi dos décadas, examinando qué grupos de trabajadores se beneficiaron más, cuáles quedaron rezagados y qué tendencias estructurales marcaron el período.

La motivación es sencilla: entender si los colombianos que trabajan hoy ganan más en términos reales que hace 17 años, y si esas ganancias fueron distribuidas de forma equitativa entre distintos grupos de la población. Los datos provienen de encuestas de hogares y están expresados en pesos constantes de 2025, lo que permite comparaciones reales descontando el efecto de la inflación.